\section{Recursive Addressing in Measurable Systems}
\label{sec:measurable}

Let $(X,\mathcal{B},T)$ be a measurable dynamical system
\cite{WaltersErgodicTheory}.  Fix a finite
family of measurable sets
\[
\mathcal{C}^{(L)} = \{C_1^{(L)},\dots,C_{r_L}^{(L)}\}\subset \mathcal{B},
\]
thought of as the observable concepts available at level $L$.  Their binary
readout process is
\[
\mathbf{b}^{(L)}_t(x)
:=
\bigl(
\ind_{C_1^{(L)}}(T^t x),\dots,\ind_{C_{r_L}^{(L)}}(T^t x)
\bigr)\in\{0,1\}^{r_L}.
\]
For a word
\[
a=(a_0,\dots,a_{\ell-1})\in \bigl(\{0,1\}^{r_L}\bigr)^\ell
\]
and an integer $t$, define the derived cylinder event
\[
C_{a,t}^{(L+1)}
:=
\{x\in X:\mathbf{b}^{(L)}_t(x)=a_0,\dots,\mathbf{b}^{(L)}_{t+\ell-1}(x)=a_{\ell-1}\}.
\]
Thus the next level is produced by finite windows in the previous one.  The
statements recorded below are classical consequences of the definitions; we
quote them for completeness because they anchor the later symbolic language.

\begin{proposition}[Recursive addressing is measurable refinement]
\label{prop:recursive-addressing-refinement}
Let
\[
\mathcal{G}^{(L)}
:=
\sigma\bigl(T^{-t}C_i^{(L)}:\ 1\le i\le r_L,\ t\in\ZZ\bigr)
\]
be the visible $\sigma$-algebra generated by the level-$L$ concepts and their
time shifts.  Then every derived cylinder event $C_{a,t}^{(L+1)}$ belongs to
$\mathcal{G}^{(L)}$.
\end{proposition}

\begin{proof}
For fixed $i$ and $t$, the map
\[
x\longmapsto \ind_{C_i^{(L)}}(T^t x)=\ind_{T^{-t}C_i^{(L)}}(x)
\]
is $\mathcal{G}^{(L)}$-measurable.  Hence the event
\[
\{x:\mathbf{b}^{(L)}_t(x)=a_0,\dots,\mathbf{b}^{(L)}_{t+\ell-1}(x)=a_{\ell-1}\}
\]
is a finite intersection of sets of the form $T^{-s}C_i^{(L)}$ and their
complements, and therefore lies in $\mathcal{G}^{(L)}$.
\end{proof}

\begin{corollary}[Visible $\sigma$-algebras do not expand]
\label{cor:recursive-addressing-visible-sigma-nonexpansion}
Suppose $\mathcal{C}^{(L+1)}$ is chosen from the derived cylinder events
$C_{a,t}^{(L+1)}$, and define
\[
\mathcal{G}^{(L+1)}
:=
\sigma\bigl(T^{-t}C:\ C\in \mathcal{C}^{(L+1)},\ t\in\ZZ\bigr).
\]
Then
\[
\mathcal{G}^{(L+1)}\subseteq \mathcal{G}^{(L)}.
\]
\end{corollary}

\begin{proof}
Each set $C\in\mathcal{C}^{(L+1)}$ is one of the derived cylinder events
$C_{a,t_0}^{(L+1)}$, so Proposition~\ref{prop:recursive-addressing-refinement}
gives $C\in\mathcal{G}^{(L)}$.  Since $\mathcal{G}^{(L)}$ is generated by all
time shifts $T^{-s}C_i^{(L)}$, it is stable under every shift $T^{-t}$:
\[
T^{-t}\bigl(T^{-s}C_i^{(L)}\bigr)=T^{-(s+t)}C_i^{(L)}\in \mathcal{G}^{(L)}.
\]
Hence $T^{-t}C\in\mathcal{G}^{(L)}$ for every $C\in\mathcal{C}^{(L+1)}$ and
every $t\in\ZZ$.  Therefore every generator of $\mathcal{G}^{(L+1)}$ lies in
$\mathcal{G}^{(L)}$, so
\[
\mathcal{G}^{(L+1)}\subseteq \mathcal{G}^{(L)}.
\]
\end{proof}

\begin{remark}
This is the measurable content of recursive addressing.  New addresses are
produced by reorganizing visible cylinder events; they do not enlarge the
observable $\sigma$-algebra.  The rest of the paper repackages this fact in
site-theoretic language.
\end{remark}

To connect the measurable setup with the symbolic site used below, fix the
alphabet
\[
\Sigma_L:=\{0,1\}^{r_L}
\]
and define the coding map
\[
\kappa_L\colon X\to \Sigma_L^{\ZZ},
\qquad
(\kappa_L(x))_t:=\mathbf{b}^{(L)}_t(x).
\]
For a finite word $a=(a_0,\dots,a_{\ell-1})\in \Sigma_L^\ell$ and an integer
$t$, write
\[
[a]_t
:=
\{y\in \Sigma_L^{\ZZ}: y_t=a_0,\dots,y_{t+\ell-1}=a_{\ell-1}\}.
\]
Let $\kappa_L^+\colon X\to \Sigma_L^{\NN}$ be the forward coding map
\[
(\kappa_L^+(x))_t:=\mathbf{b}^{(L)}_t(x)
\qquad (t\ge 0),
\]
and for $a\in \Sigma_L^{<\infty}$ let $C_a\subseteq \Sigma_L^{\NN}$ denote the
usual one-sided cylinder of forward sequences beginning with $a$.  Define the
admissible forward words by
\[
\Aadm^{(L)}
:=
\{a\in \Sigma_L^{<\infty}: (\kappa_L^+)^{-1}(C_a)\neq\varnothing\}.
\]

\begin{proposition}[Coding bridge from measurable events to prefix cylinders]
\label{prop:measurable-coding-prefix-bridge}
For every finite word $a=(a_0,\dots,a_{\ell-1})\in \Sigma_L^\ell$ and every
$t\in\ZZ$,
\[
C_{a,t}^{(L+1)}=\kappa_L^{-1}([a]_t).
\]
In particular, for every $a\in \Aadm^{(L)}$,
\[
(\kappa_L^+)^{-1}(C_a)=C_{a,0}^{(L+1)},
\]
so the forward symbolic cylinders are exactly the measurable derived cylinder
events at time $0$.  More generally,
\[
\kappa_L^{-1}(\sigma^{-t}[a]_0)=C_{a,t}^{(L+1)},
\]
where $\sigma$ is the shift on $\Sigma_L^{\ZZ}$.
\end{proposition}

\begin{proof}
By definition, $x\in C_{a,t}^{(L+1)}$ if and only if
\[
\mathbf{b}^{(L)}_t(x)=a_0,\dots,\mathbf{b}^{(L)}_{t+\ell-1}(x)=a_{\ell-1},
\]
which is equivalent to $\kappa_L(x)\in [a]_t$.  This proves the first
identity.  The statement for $t=0$ is exactly the definition of the one-sided
cylinder $C_a$, and the shifted form is the same identity written as a pullback
of a translated cylinder in $\Sigma_L^{\ZZ}$.
\end{proof}

\section{Prefix Sites and Readout Presheaves}
\label{sec:prefix-sites}

From now on we stay inside the admissible image shift
\[
\Xadm:=\im(\kappa_L^+)\subseteq \Sigma_L^{\NN}
\]
produced by the measurable coding map, so every admissible word is literally a
trace that occurs in $\Xadm$.
Proposition~\ref{prop:measurable-coding-prefix-bridge} shows that, after fixing
a time origin, the symbolic site is obtained from the measurable readout
process by pulling back cylinders inside $\Xadm$, not inside the full shift,
so every object considered below is a subset of $\Xadm$.

\subsection{The Cylinder Category}

Fix a finite alphabet $\Sigma$ and let $\Xadm\subseteq \Sigma^{\NN}$ be any
subshift that arises as an admissible image $\im(\kappa_L^+)$ of a measurable
coding map.  Let $\Sigma^{<\infty}$ be the set of finite words, write
$a\preceq b$ when $a$ is a prefix of $b$, and set
\[
\Aadm:=\{a\in\Sigma^{<\infty}:\ C_a^{\mathrm{full}}\cap \Xadm\neq\varnothing\},
\]
the set of words that actually occur inside $\Xadm$.  For such a word $a$, let
$C_a\subseteq \Xadm$ be the trace cylinder consisting of all sequences of
$\Xadm$ that begin with $a$
\cite{LindMarcus1995SymbolicDynamics,Kitchens1998SymbolicDynamics}; equivalently,
$C_a=\Xadm\cap C_a^{\mathrm{full}}$, where $C_a^{\mathrm{full}}$ is the usual
cylinder in the full shift $\Sigma^{\NN}$.

\begin{definition}[Cylinder algebra]
\label{def:cylinder-algebra}
The \emph{cylinder algebra} $\mathfrak{C}(\Aadm)$ is the Boolean subalgebra of
$\mathcal{P}(\Xadm)$ generated by the trace cylinders $\{C_a:a\in\Aadm\}$.  Its
elements are finite unions of admissible cylinders.  Because $\Xadm$ consists
only of actual traces, the following properties hold automatically.
\begin{enumerate}[label=\textup{(CA\arabic*)}]
\item\label{ax:nesting} if $b\succeq a$, then $C_b\subseteq C_a$;
\item\label{ax:intersection} if $a,b\in\Aadm$ and $C_a\cap C_b\neq\varnothing$,
then there exists $d\in\Aadm$ with $d\succeq a$, $d\succeq b$, and
$C_d = C_a\cap C_b$ (take $d$ to be the longer of the two words);
\item\label{ax:refinement} every nonempty element of $\mathfrak{C}(\Aadm)$ can
be written as a finite disjoint union of cylinders $C_{a_i}$ with
$a_i\in\Aadm$.
\end{enumerate}
\end{definition}

Axiom~\ref{ax:intersection} simply records that if a sequence in $\Xadm$
begins with both $a$ and $b$, then one word extends the other, so the overlap
agrees with the cylinder of the longer word.  The refinement axiom follows
because every finite union of basic trace cylinders can be refined by
splitting into disjoint cylinders of a common depth.  When $\Xadm$ is the full
shift, these statements reduce to the usual combinatorics of prefixes in
$\Sigma^{\NN}$.

\begin{definition}[Prefix site]
\label{def:prefix-site}
The \emph{prefix site} $\mathsf{Pref}(\Aadm)$ is the category whose objects
are the elements of the cylinder algebra $\mathfrak{C}(\Aadm)$ and whose
morphisms are inclusions.  A finite family $\{U_i\hookrightarrow U\}_{i=1}^n$
is declared to be a \emph{cover} if
\[
U = \bigcup_{i=1}^n U_i.
\]
\end{definition}

When convenient, we identify the cylinder $C_a$ with its address $a$ and write
$b\to a$ for the inclusion $C_b\hookrightarrow C_a$ when $b\succeq a$.  In
this notation, a \emph{cylinder-algebra cover} of $a$ is a finite family
$\{b_i\to a\}_{i=1}^n$ such that $C_a=\bigcup_i C_{b_i}$.

\begin{proposition}[Cylinder-algebra covers generate a Grothendieck topology]
\label{prop:prefix-site-grothendieck-topology}
The cylinder-algebra covers satisfy the identity, base-change, and transitivity axioms
of a finitary pretopology on $\mathsf{Pref}(\Aadm)$
\cite[Chapter~III]{MacLaneMoerdijk1992},
\cite[\S C2.1]{Johnstone2002Elephant}.  Hence they generate a
Grothendieck topology $J_{\mathrm{cyl}}$, and
\[
(\mathsf{Pref}(\Aadm),J_{\mathrm{cyl}})
\]
is a well-defined site.
\end{proposition}

\begin{proof}
The singleton family $\{U\hookrightarrow U\}$ is a cover because $U=U$.

For base change, let $\{U_i\hookrightarrow U\}_{i=1}^n$ be a cover and let
$V\hookrightarrow U$ be a morphism (inclusion).  Then
\[
V
=
V\cap U
=
\bigcup_{i=1}^n (V\cap U_i).
\]
Each $V\cap U_i$ belongs to $\mathfrak{C}(\Aadm)$ by
axiom~\ref{ax:intersection} (applied iteratively to the cylinder
decompositions of $V$ and $U_i$) and axiom~\ref{ax:refinement}.  Writing each
nonempty $V\cap U_i$ as a finite union of admissible cylinders gives a cover of
$V$ that refines the original cover of $U$.

For transitivity, if $\{U_i\hookrightarrow U\}_{i=1}^n$ is a cover and each
$\{V_{ij}\hookrightarrow U_i\}_{j=1}^{m_i}$ is a cover of $U_i$, then
\[
U
=
\bigcup_{i=1}^n U_i
=
\bigcup_{i=1}^n \bigcup_{j=1}^{m_i} V_{ij},
\]
so the composite family is again a cover.
Therefore the covers form a finitary pretopology and hence generate a
Grothendieck topology \cite[Proposition~III.2.2]{MacLaneMoerdijk1992}.
\end{proof}

\subsection{The Readout Prestack}

The paper keeps the readout object deliberately abstract.  A local readout over
$a$ is any object that can be evaluated using only the information visible on
$C_a$ and that remains admissible under further refinement.

\begin{definition}[Readout prestack]
\label{def:prefix-sheaf-readout}
A \emph{readout prestack} on $\mathsf{Pref}(\Aadm)$ is a contravariant pseudofunctor
\[
\Rpre:\mathsf{Pref}(\Aadm)^{\mathrm{op}}\to \mathbf{Gpd}.
\]
For $a\in\Aadm$, the groupoid $\Rpre(a)$ records the admissible local readout
objects over $C_a$ and their internal morphisms.  For a refinement $b\succeq a$,
the restriction functor
\[
\rho_{a\leftarrow b}\colon \Rpre(a)\longrightarrow \Rpre(b)
\]
is functorial and encodes admissibility under further refinement.  We write
$\Ob\Rpre(a)$ for its object set and $\Aut_{\Rpre}(s)$ for the automorphism
group of $s\in\Ob\Rpre(a)$.  Whenever we state that $\Rpre(a)=\varnothing$, we
mean that $\Ob\Rpre(a)$ has no objects, and all previous set-valued terminology
applies to $\Ob\Rpre$ via the forgetful functor $\mathbf{Gpd}\to\mathbf{Set}$.
We say that $\Rpre$ is \emph{separated} with respect to a cover
$\{b_i\to a\}_{i=1}^n$ if the canonical map
$\Ob\Rpre(a)\to\prod_i\Ob\Rpre(b_i)$ is injective, equivalently if the
underlying set-valued presheaf of objects is separated on that cover
\cite[\S II.1]{MacLaneMoerdijk1992}.
\end{definition}

\section{Local Sections and Descent}
\label{sec:descent}
\label{sec:descent-groupoid}

\subsection{Local Sections}

\begin{proposition}[Addressability as existence of a local section]
\label{prop:null-as-local-section-obstruction}
For an admissible address $a\in\Aadm$, readout is internally defined at $a$ if
and only if
\[
\Rpre(a)\neq\varnothing.
\]
\end{proposition}

\begin{proof}
By Definition~\ref{def:prefix-sheaf-readout}, the object set
$\Ob\Rpre(a)$ consists exactly of the admissible local readout objects over
the cylinder $C_a$.  Therefore $\Rpre(a)\neq\varnothing$ (meaning
$\Ob\Rpre(a)\neq\varnothing$) if and only if at least one such local readout
exists.
\end{proof}

\begin{corollary}[Three standard failure modes]
\label{cor:null-trichotomy-local-section}
For an admissible finite-level address $a\in\Aadm$, failure of a well-defined
readout at $a$ has three internal
forms:
\begin{enumerate}[label=\textup{(\arabic*)}]
\item \emph{empty local fiber}: $\Rpre(a)=\varnothing$ (i.e., $\Ob\Rpre(a)$ has
no objects);
\item \emph{descent failure}: local sections exist on a cylinder-algebra cover of $a$
and are pairwise compatible, but no object of $\Rpre(a)$ restricts to them;
\item \emph{non-separatedness}: two distinct objects of $\Rpre(a)$ have the
same restrictions to a cylinder-algebra cover of $a$.
\end{enumerate}
These modes are not mutually exclusive, and the latter two are relative to a
chosen cylinder-algebra cover of $a$.
\end{corollary}

\begin{proof}
The first case is Proposition~\ref{prop:null-as-local-section-obstruction}.  If
$\Rpre(a)\neq\varnothing$ but a compatible family on a cylinder-algebra cover does not
come from any object of $\Rpre(a)$, then the failure is one of effective
descent.  If distinct objects of $\Rpre(a)$ induce the same compatible family
on a cover, then the failure is one of separatedness.  These are the three
standard internal failure modes for the existence and uniqueness of readout on
a presheaf over a fixed admissible object and a fixed cover
\cite[\S II.1--2]{MacLaneMoerdijk1992}.
\end{proof}

\begin{example}[Binary prefix site and the three failure modes]
\label{ex:binary-prefix-site-failure-modes}
Let $X=\{0,1\}^{\NN}$, let $\Sigma=\{0,1\}$, let $\Aadm=\Sigma^{<\infty}$, and
let $C_a\subseteq X$ be the usual cylinder of sequences beginning with the word
$a$.  The root address $\varepsilon$ is covered by the two depth-one cylinders
$0\to \varepsilon$ and $1\to \varepsilon$.

Three explicit readout presheaves on this same binary prefix site realize the
three failure modes from Corollary~\ref{cor:null-trichotomy-local-section}.
\begin{enumerate}[label=\textup{(\arabic*)}]
\item \emph{Empty fiber}: let $\Rpre_{\mathrm{emp}}(\varepsilon)=\varnothing$
and $\Rpre_{\mathrm{emp}}(b)=\{*\}$ for every nonempty word $b$, with the
unique possible restriction maps.  Readout fails immediately at the root.
\item \emph{Descent failure}: let $\Rpre_{\mathrm{desc}}(\varepsilon)=\{g\}$,
let $\Rpre_{\mathrm{desc}}(0)=\{x_0,y_0\}$ and
$\Rpre_{\mathrm{desc}}(1)=\{x_1,y_1\}$, and let every deeper cylinder below
$0$ or $1$ inherit the corresponding two-point set with identity restrictions.
The two restriction maps from $\varepsilon$ send $g$ to $x_0$ and $x_1$.
The pair overlap $C_0\cap C_1$ is empty, so every pair of local sections on the
cover is automatically compatible.  In particular, the family $(y_0,y_1)$ is
compatible on the cover
$\{0,1\to\varepsilon\}$ but does not descend to a global section.
\item \emph{Non-separatedness}: let
$\Rpre_{\mathrm{ns}}(\varepsilon)=\{u,v\}$ and
$\Rpre_{\mathrm{ns}}(b)=\{*\}$ for every nonempty word $b$, with both
restrictions from $\varepsilon$ constant.  Then $u\neq v$, but their
restrictions to the cover $\{0,1\to\varepsilon\}$ agree.
\end{enumerate}
This concrete binary site keeps the geometry fixed while isolating the three
internal presheaf-theoretic mechanisms.
\end{example}

\subsection{Intrinsic Comparison Torsors}
\label{subsec:comparison-torsors}

Let $\mathfrak{U}=\{b_i\to a\}_{i=1}^n$ be a finite cylinder-algebra cover.  Because
the cylinder algebra is closed under finite intersections
(Definition~\ref{def:cylinder-algebra}), every nonempty overlap
$C_{b_{i_0}}\cap\cdots\cap C_{b_{i_q}}$ is itself a cylinder
$C_{b_{i_0\dots i_q}}$ refining each $b_{i_r}$.  For an object
$s_i\in\Ob\Rpre(b_i)$ and such an overlap we abbreviate
\[
s_{i|i_0\dots i_q}
:=
\rho_{b_i\leftarrow b_{i_0\dots i_q}}(s_i)\in\Ob\Rpre(b_{i_0\dots i_q}).
\]
In particular, for a pair overlap $b_{ij}$ we obtain the comparison groupoid
\[
\Isom_{ij}:=
\Isom_{\Rpre(b_{ij})}\bigl(s_{j|ij},s_{i|ij}\bigr),
\]
and when this set is nonempty it is acted on simply transitively by the
automorphism group of the target object $s_{i|ij}$ via postcomposition.

\begin{definition}[Abelian band along a cover]\label{def:abelian-band}
Fix local objects $\{s_i\in\Ob\Rpre(b_i)\}_{i=1}^n$.  An \emph{abelian band}
for $\{s_i\}$ consists of an abelian group $A$ together with isomorphisms
\[
\vartheta_{i,V}\colon \Aut_{\Rpre}(s_{i|V})\xrightarrow{\;\cong\;} A
\]
for every nonempty refinement $V\to b_i$ obtained from a finite intersection
of members of $\mathfrak{U}$, such that $\vartheta_{i,V}$ is natural under
further restriction: if $V'\to V$ is induced by an additional intersection,
then $\vartheta_{i,V'}\circ \rho_{V\leftarrow V'}=\vartheta_{i,V}$.  Equivalently,
the family $\{\Aut_{\Rpre}(s_{i|V})\}$ is identified with the constant band $A$
along the cover.
\end{definition}

Under this hypothesis every comparison set inherits a canonical $A$-torsor
structure: for $\xi\in \Isom_{ij}$ and $a\in A$ we set
\[
a\cdot \xi := \vartheta_{i,b_{ij}}^{-1}(a)\circ \xi.
\]
This action is free and transitive, and it is compatible with composition on
triple overlaps because the identifications $\vartheta_{i,V}$ are natural.
Consequently the torsors and their composition maps required for the defect
construction are intrinsic to $\Rpre$ as soon as an abelian band has been
specified; no extra comparison datum has to be prescribed externally.
Moreover, for all $a,b\in A$ and compatible comparison elements we have
\begin{equation}\label{eq:torsor-additivity}
m_{ijk}(a\cdot \xi,b\cdot \eta) = (a+b)\cdot m_{ijk}(\xi,\eta),
\end{equation}
since the action is given by postcomposition with automorphisms of the target.

\subsection{The \v{C}ech Complex of a Cylinder-Algebra Cover}
\label{subsec:cech-complex}

The groupoid-valued prestack $\Rpre$ therefore produces canonical torsors of
comparisons on every overlap once an abelian band is fixed.  To speak about a
\v{C}ech defect, we now review the underlying cochain complex.  The general
framework is classical; see \cite[Tag~01FG]{StacksProject} for the abstract
\v{C}ech complex, \cite[Chapter~V]{Bredon1997SheafTheory} for
sheaf-cohomological background, and
\cite[\S IV.3]{Giraud1971CohomologieNonAbelienne} for the banded (possibly
non-abelian) refinement.

\begin{lemma}[Normalized \v{C}ech complex of a finite cylinder-algebra cover]
\label{lem:cech-complex-formal}
Let $\mathfrak{U}=\{b_i\to a\}_{i=1}^n$ be a finite cylinder-algebra cover of
$a\in\Aadm$, and let $A$ be an abelian group.  Define the \emph{ordered cover
nerve} $N(\mathfrak{U})$ as the abstract simplicial complex on
$\{1,\dots,n\}$ whose $q$-simplices are the strictly increasing tuples
\[
N_q(\mathfrak{U})
:=
\left\{
(i_0,\dots,i_q): 1\le i_0<\cdots<i_q\le n,\ 
C_{b_{i_0}}\cap\cdots\cap C_{b_{i_q}}\neq\varnothing
\right\}.
\]
Define the normalized \v{C}ech cochain groups
\[
\check{C}^q(\mathfrak{U},A)
:=
\prod_{(i_0,\dots,i_q)\in N_q(\mathfrak{U})} A.
\]
The coboundary $\delta^q\colon \check{C}^q\to \check{C}^{q+1}$ is given by
\[
(\delta^q f)_{i_0\dots i_{q+1}}
:=
\sum_{s=0}^{q+1} (-1)^s\, f_{i_0\dots\widehat{i_s}\dots i_{q+1}},
\]
for $i_0<\cdots<i_{q+1}$, where $\widehat{i_s}$ denotes omission.  Then
$\delta^{q+1}\circ\delta^q=0$, and
\[
\check{H}^q(\mathfrak{U},A)
:=
\ker\delta^q / \im\delta^{q-1}
\]
is the $q$-th normalized \v{C}ech cohomology group of $\mathfrak{U}$ with
coefficients in $A$.  In particular, a $2$-cocycle is a family
$(\alpha_{ijk})_{i<j<k}\in\check{C}^2(\mathfrak{U},A)$ satisfying
\begin{equation}\label{eq:cocycle-condition}
\alpha_{jkl} - \alpha_{ikl} + \alpha_{ijl} - \alpha_{ijk} = 0
\end{equation}
for every ordered quadruple $i<j<k<l$ with
$C_{b_i}\cap C_{b_j}\cap C_{b_k}\cap C_{b_l}\neq\varnothing$, and two
$2$-cocycles $\alpha,\alpha'$ are
cohomologous when $\alpha-\alpha'=\delta^1 h$ for some
$1$-cochain $(h_{ij})_{i<j}\in\check{C}^1(\mathfrak{U},A)$.
\end{lemma}

\begin{proof}
This is the standard simplicial or normalized \v{C}ech complex of an ordered
finite nerve.  The identity $\delta^{q+1}\circ\delta^q=0$ is the usual
alternating-sum cancellation, and the remaining assertions are definitions.
See \cite{HatcherAlgebraicTopology} for simplicial cochains and
\cite[Tag~01FG]{StacksProject} for the \v{C}ech formulation.
\end{proof}

\subsection{\v{C}ech-Type Gluing Obstructions}

\begin{definition}[Intrinsic $A$-banded comparison datum and \v{C}ech defect]
\label{def:prefix-site-h2-gluing-obstruction}
Fix an admissible address $a\in\Aadm$, a cylinder-algebra cover
$\mathfrak{U}=\{b_i\to a\}_{i=1}^n$, and an abelian group $A$.  Suppose
$\{s_i\in\Ob\Rpre(b_i)\}$ is a family of local sections equipped with an
abelian band in the sense of Definition~\ref{def:abelian-band}.  For every
nonempty pair overlap $b_{ij}$ we form the comparison torsor
\[
\Comp_{ij}
:=
\Isom_{\Rpre(b_{ij})}(s_{j|ij},s_{i|ij}),
\]
viewed as an $A$-torsor via the action $a\cdot\xi
:=\vartheta_{i,b_{ij}}^{-1}(a)\circ\xi$.  On each nonempty ordered triple
overlap the composition map
\[
m_{ijk}\colon \Comp_{ij}\times\Comp_{jk}\longrightarrow \Comp_{ik},
\qquad
m_{ijk}(\xi,\eta):=\xi\circ\eta,
\]
is $A$-equivariant and associative by construction
\cite[\S IV.3]{Giraud1971CohomologieNonAbelienne}.  Choosing elements
$\gamma_{ij}\in \Comp_{ij}$ for all nonempty ordered pair overlaps therefore
produces the data previously listed.

Because $\Comp_{ik}$ is an $A$-torsor, on each nonempty ordered triple overlap
there is a unique element $\alpha_{ijk}\in A$ such that
\begin{equation}\label{eq:defect-via-composition}
m_{ijk}(\gamma_{ij},\gamma_{jk})=\alpha_{ijk}\cdot \gamma_{ik}.
\end{equation}
The family $(\alpha_{ijk})$ is the associated \emph{\v{C}ech defect}.  The
associativity of the composition maps guarantees that $(\alpha_{ijk})$
satisfies the cocycle condition~\eqref{eq:cocycle-condition} on every
$4$-fold overlap, so it defines a class
$[\alpha]\in\check{H}^2(\mathfrak{U},A)$.
\end{definition}

The next lemma records the effect of changing representatives.

\begin{lemma}[Change of trivialization]\label{lem:change-trivialization}
Let $\beta=(\beta_{ij})_{i<j}\in\check{C}^1(\mathfrak{U},A)$ and set
$\gamma'_{ij}:=\beta_{ij}\cdot \gamma_{ij}$.  Then the defect computed from the
family $\{\gamma'_{ij}\}$ is the cohomologous cocycle
$\alpha'=\alpha+\delta^1\beta$.
\end{lemma}

\begin{proof}
Equation~\eqref{eq:torsor-additivity} gives
\[
m_{ijk}(\gamma'_{ij},\gamma'_{jk})
=
m_{ijk}(\beta_{ij}\cdot\gamma_{ij},\beta_{jk}\cdot \gamma_{jk})
=
(\beta_{ij}+\beta_{jk})\cdot m_{ijk}(\gamma_{ij},\gamma_{jk}).
\]
Substituting~\eqref{eq:defect-via-composition} produces
\[
m_{ijk}(\gamma'_{ij},\gamma'_{jk})
=
(\beta_{ij}+\beta_{jk}+\alpha_{ijk})\cdot \gamma_{ik}
=
(\beta_{ij}+\beta_{jk}+\alpha_{ijk}-\beta_{ik})\cdot \gamma'_{ik},
\]
which is precisely the statement that $\alpha'=\alpha+\delta^1\beta$.  Hence
the class $[\alpha]$ depends only on the banded comparison datum.
\end{proof}

\begin{proposition}[Nontrivial \v{C}ech defect prevents coherent global readout]
\label{prop:null-as-h2-obstruction}
In the setting of Definition~\ref{def:prefix-site-h2-gluing-obstruction},
suppose further that there exists a global section $s\in\Ob\Rpre(a)$ whose
restrictions to the cover are the sections $\{s_i\}$ and whose induced
comparisons on every pair overlap $(i,j)$ produce distinguished elements
$\gamma_{ij}^s\in\Comp_{ij}$ that are compatible with the composition maps:
\[
m_{ijk}(\gamma_{ij}^s,\gamma_{jk}^s)=\gamma_{ik}^s
\]
on every triple overlap.  Then the defect class $[\alpha]$ computed from the
induced elements $\gamma^s_{ij}$ is trivial.  In particular, if
$[\alpha]\neq 0$, then no such global section $s$ exists.
\end{proposition}

\begin{proof}
The global section $s$ restricts to $s_i$ on each $b_i$, and on every pair
overlap the induced comparison between $s_j$ and $s_i$ via $s$ gives a
distinguished element $\gamma_{ij}^s$.  On every triple overlap these compose
coherently:
\[
m_{ijk}(\gamma_{ij}^s,\gamma_{jk}^s)=\gamma_{ik}^s.
\]
By Definition~\ref{def:prefix-site-h2-gluing-obstruction}, this means
$\alpha_{ijk}=0$ on every triple overlap.  Hence the resulting defect cocycle
is zero, so $[\alpha]=0$.
\end{proof}

